% Options for packages loaded elsewhere
\PassOptionsToPackage{unicode}{hyperref}
\PassOptionsToPackage{hyphens}{url}
\documentclass[
]{article}
\usepackage{xcolor}
\usepackage[margin=1in]{geometry}
\usepackage{amsmath,amssymb}
\setcounter{secnumdepth}{-\maxdimen} % remove section numbering
\usepackage{iftex}
\ifPDFTeX
  \usepackage[T1]{fontenc}
  \usepackage[utf8]{inputenc}
  \usepackage{textcomp} % provide euro and other symbols
\else % if luatex or xetex
  \usepackage{unicode-math} % this also loads fontspec
  \defaultfontfeatures{Scale=MatchLowercase}
  \defaultfontfeatures[\rmfamily]{Ligatures=TeX,Scale=1}
\fi
\usepackage{lmodern}
\ifPDFTeX\else
  % xetex/luatex font selection
\fi
% Use upquote if available, for straight quotes in verbatim environments
\IfFileExists{upquote.sty}{\usepackage{upquote}}{}
\IfFileExists{microtype.sty}{% use microtype if available
  \usepackage[]{microtype}
  \UseMicrotypeSet[protrusion]{basicmath} % disable protrusion for tt fonts
}{}
\makeatletter
\@ifundefined{KOMAClassName}{% if non-KOMA class
  \IfFileExists{parskip.sty}{%
    \usepackage{parskip}
  }{% else
    \setlength{\parindent}{0pt}
    \setlength{\parskip}{6pt plus 2pt minus 1pt}}
}{% if KOMA class
  \KOMAoptions{parskip=half}}
\makeatother
\usepackage{color}
\usepackage{fancyvrb}
\newcommand{\VerbBar}{|}
\newcommand{\VERB}{\Verb[commandchars=\\\{\}]}
\DefineVerbatimEnvironment{Highlighting}{Verbatim}{commandchars=\\\{\}}
% Add ',fontsize=\small' for more characters per line
\usepackage{framed}
\definecolor{shadecolor}{RGB}{248,248,248}
\newenvironment{Shaded}{\begin{snugshade}}{\end{snugshade}}
\newcommand{\AlertTok}[1]{\textcolor[rgb]{0.94,0.16,0.16}{#1}}
\newcommand{\AnnotationTok}[1]{\textcolor[rgb]{0.56,0.35,0.01}{\textbf{\textit{#1}}}}
\newcommand{\AttributeTok}[1]{\textcolor[rgb]{0.13,0.29,0.53}{#1}}
\newcommand{\BaseNTok}[1]{\textcolor[rgb]{0.00,0.00,0.81}{#1}}
\newcommand{\BuiltInTok}[1]{#1}
\newcommand{\CharTok}[1]{\textcolor[rgb]{0.31,0.60,0.02}{#1}}
\newcommand{\CommentTok}[1]{\textcolor[rgb]{0.56,0.35,0.01}{\textit{#1}}}
\newcommand{\CommentVarTok}[1]{\textcolor[rgb]{0.56,0.35,0.01}{\textbf{\textit{#1}}}}
\newcommand{\ConstantTok}[1]{\textcolor[rgb]{0.56,0.35,0.01}{#1}}
\newcommand{\ControlFlowTok}[1]{\textcolor[rgb]{0.13,0.29,0.53}{\textbf{#1}}}
\newcommand{\DataTypeTok}[1]{\textcolor[rgb]{0.13,0.29,0.53}{#1}}
\newcommand{\DecValTok}[1]{\textcolor[rgb]{0.00,0.00,0.81}{#1}}
\newcommand{\DocumentationTok}[1]{\textcolor[rgb]{0.56,0.35,0.01}{\textbf{\textit{#1}}}}
\newcommand{\ErrorTok}[1]{\textcolor[rgb]{0.64,0.00,0.00}{\textbf{#1}}}
\newcommand{\ExtensionTok}[1]{#1}
\newcommand{\FloatTok}[1]{\textcolor[rgb]{0.00,0.00,0.81}{#1}}
\newcommand{\FunctionTok}[1]{\textcolor[rgb]{0.13,0.29,0.53}{\textbf{#1}}}
\newcommand{\ImportTok}[1]{#1}
\newcommand{\InformationTok}[1]{\textcolor[rgb]{0.56,0.35,0.01}{\textbf{\textit{#1}}}}
\newcommand{\KeywordTok}[1]{\textcolor[rgb]{0.13,0.29,0.53}{\textbf{#1}}}
\newcommand{\NormalTok}[1]{#1}
\newcommand{\OperatorTok}[1]{\textcolor[rgb]{0.81,0.36,0.00}{\textbf{#1}}}
\newcommand{\OtherTok}[1]{\textcolor[rgb]{0.56,0.35,0.01}{#1}}
\newcommand{\PreprocessorTok}[1]{\textcolor[rgb]{0.56,0.35,0.01}{\textit{#1}}}
\newcommand{\RegionMarkerTok}[1]{#1}
\newcommand{\SpecialCharTok}[1]{\textcolor[rgb]{0.81,0.36,0.00}{\textbf{#1}}}
\newcommand{\SpecialStringTok}[1]{\textcolor[rgb]{0.31,0.60,0.02}{#1}}
\newcommand{\StringTok}[1]{\textcolor[rgb]{0.31,0.60,0.02}{#1}}
\newcommand{\VariableTok}[1]{\textcolor[rgb]{0.00,0.00,0.00}{#1}}
\newcommand{\VerbatimStringTok}[1]{\textcolor[rgb]{0.31,0.60,0.02}{#1}}
\newcommand{\WarningTok}[1]{\textcolor[rgb]{0.56,0.35,0.01}{\textbf{\textit{#1}}}}
\usepackage{graphicx}
\makeatletter
\newsavebox\pandoc@box
\newcommand*\pandocbounded[1]{% scales image to fit in text height/width
  \sbox\pandoc@box{#1}%
  \Gscale@div\@tempa{\textheight}{\dimexpr\ht\pandoc@box+\dp\pandoc@box\relax}%
  \Gscale@div\@tempb{\linewidth}{\wd\pandoc@box}%
  \ifdim\@tempb\p@<\@tempa\p@\let\@tempa\@tempb\fi% select the smaller of both
  \ifdim\@tempa\p@<\p@\scalebox{\@tempa}{\usebox\pandoc@box}%
  \else\usebox{\pandoc@box}%
  \fi%
}
% Set default figure placement to htbp
\def\fps@figure{htbp}
\makeatother
\setlength{\emergencystretch}{3em} % prevent overfull lines
\providecommand{\tightlist}{%
  \setlength{\itemsep}{0pt}\setlength{\parskip}{0pt}}
\usepackage{bookmark}
\IfFileExists{xurl.sty}{\usepackage{xurl}}{} % add URL line breaks if available
\urlstyle{same}
\hypersetup{
  pdftitle={Effect of Pyrolysis condition on biochar},
  pdfauthor={Sodiq Yusuf},
  hidelinks,
  pdfcreator={LaTeX via pandoc}}

\title{Effect of Pyrolysis condition on biochar}
\author{Sodiq Yusuf}
\date{2025-12-08}

\begin{document}
\maketitle

\begin{verbatim}
## Warning: package 'DBI' was built under R version 4.4.3
\end{verbatim}

\begin{verbatim}
## Warning: package 'RSQLite' was built under R version 4.4.3
\end{verbatim}

\begin{verbatim}
## Warning: package 'tidyverse' was built under R version 4.4.3
\end{verbatim}

\begin{verbatim}
## Warning: package 'ggplot2' was built under R version 4.4.3
\end{verbatim}

\begin{verbatim}
## Warning: package 'tidyr' was built under R version 4.4.3
\end{verbatim}

\begin{verbatim}
## Warning: package 'readr' was built under R version 4.4.3
\end{verbatim}

\begin{verbatim}
## Warning: package 'purrr' was built under R version 4.4.3
\end{verbatim}

\begin{verbatim}
## Warning: package 'dplyr' was built under R version 4.4.3
\end{verbatim}

\begin{verbatim}
## Warning: package 'stringr' was built under R version 4.4.3
\end{verbatim}

\begin{verbatim}
## Warning: package 'forcats' was built under R version 4.4.3
\end{verbatim}

\begin{verbatim}
## Warning: package 'lubridate' was built under R version 4.4.3
\end{verbatim}

\begin{verbatim}
## -- Attaching core tidyverse packages ------------------------ tidyverse 2.0.0 --
## v dplyr     1.1.4     v readr     2.1.5
## v forcats   1.0.1     v stringr   1.5.1
## v ggplot2   4.0.0     v tibble    3.2.1
## v lubridate 1.9.4     v tidyr     1.3.1
## v purrr     1.1.0     
## -- Conflicts ------------------------------------------ tidyverse_conflicts() --
## x dplyr::filter() masks stats::filter()
## x dplyr::lag()    masks stats::lag()
## i Use the conflicted package (<http://conflicted.r-lib.org/>) to force all conflicts to become errors
\end{verbatim}

' \# Chapter 1: introduction Pyrolysis is the thermal decomposition of a
substance in the absence of air. The products of this process include
solid biochar, liquid bio-oil, and gaseous bio-gas. Recent research have
explored pyrolysis technology as a viable technique to convert wastes to
useful products such as biochar, bio-oil and gases. Biochar is the black
solid residue that have great potential for carbon sequestration,
environmental remediation, biodegradability.

The rapid advancements in Additive Manufacturing (AM), commonly known as
3D printing, have revolutionized various industries by enabling the
production of complex geometries with high precision and customization
(Vidakis, 2024) (Bolanakis, 2024). Concurrently, there is a growing
imperative to transition towards sustainable materials in manufacturing
processes, driven by environmental concerns, resource depletion, and the
need to reduce carbon footprints (Vidakis, 2024) (Bolanakis, 2025). This
demand has spurred significant interest in developing eco-friendly
composites for AM applications, utilizing renewable resources to replace
conventional petroleum-based materials (Vidakis, 2023).

Biochar, a carbon-rich material produced from the thermochemical
conversion of biomass through pyrolysis, has emerged as a promising
sustainable filler for polymer composites in AM (Bolanakis, 2024)
(Aboughaly, 2023). Its superior porosity, large specific surface area,
and high carbon content make it an attractive additive for enhancing the
functionality and performance of 3D-printed parts (Bolanakis, 2024).
Research has demonstrated that incorporating biochar into various
polymer matrices, such as high-density polyethylene (HDPE), polyethylene
terephthalate glycol (PETG), and polylactic acid (PLA), can
significantly improve mechanical properties, including tensile strength,
flexural strength, and impact resistance, while also contributing to the
eco-friendly profile of the final product (Vidakis, 2024) (Bolanakis,
2025) (Vidakis, 2023) (George, 2023) (Gkiliopoulos, 2025). For instance,
biochar has been shown to increase the tensile strength by over 20\% and
flexural strength by 35.9\% in HDPE composites at optimal loadings
(Vidakis, 2024).

The properties of biochar are not static but are highly dependent on the
pyrolysis conditions employed during its production, particularly
temperature and retention time (Wijitkosum, 2023) (Suman, 2017)
(Fernandes, 2020). These parameters dictate critical characteristics
such as carbon content, nitrogen content, carbon to nitrogen (C/N)
ratio, density, and structural features reflected by the Raman factor
(Wijitkosum, 2023) (Fernandes, 2020) (Zhang, 2024) (Venkatesh, 2022).
For example, increasing pyrolysis temperature generally leads to higher
carbon content, increased aromaticity (observable via Raman
spectroscopy), and altered density, while nitrogen content and C/N ratio
can vary, influencing properties like water-holding capacity and
interfacial interactions within a composite (Wijitkosum, 2023)
(Fernandes, 2020) (Zhang, 2024) (Venkatesh, 2022) (Hasan, 2019). These
intrinsic properties of biochar, in turn, play a crucial role in
determining its effectiveness as a reinforcing agent and its
compatibility with polymer matrices in AM (Aboughaly, 2023).

Despite the growing interest in biochar for AM, a comprehensive
understanding of the direct correlation between specific pyrolysis
conditions (temperature and retention time), the resulting biochar
properties (carbon content, nitrogen content, C/N ratio, density, and
Raman factor), and their subsequent impact on the mechanical performance
(flexural strength and flexural modulus) of additive manufactured
composites remains an area requiring further detailed investigation.
Optimizing these relationships is critical for tailoring biochar to
achieve desired material properties in AM applications, thereby
maximizing its potential as a sustainable and high-performance filler.

Therefore, this research aims to investigate the effect of pyrolysis
conditions (temperature and retention time) on the properties of biochar
(carbon content, nitrogen content, carbon to nitrogen ratio, density,
and Raman factor) and how these biochar properties, in turn, affect the
mechanical properties (flexural strength, flexural modulus) of
composites prepared from the biochar through additive manufacturing.

\subsection{1.1 Research objectives}\label{research-objectives}

This research aims to investigate the effect of pyrolysis condition
(temperature and retention time) on the properties of biochar(carbon
content, nitrogen content, carbon to nitrogen ratio,density, and raman
factor) and how this affect the mechanical properties (flexural
strength, flexural modulus) of composites prepared from the biochar
through additive manufacturing.

\subsection{1.2 Justification of the
research}\label{justification-of-the-research}

A major challenge in the utilization of biobased materials for additive
manufacturing is their hydrophobicity which subject them to decay with
time especially in outdoor applications. Biochar are carbon rich
hydrophobic material that can be produced from pyrolysis of biobased
materials especially forest residues and wastes. In this research, Wood
wastes are converted to biochar and the effect of the pyrolysis
condition on biochar for additive manufacturing applications was
investigated.

\subsection{1.3 References}\label{references}

\begin{itemize}
\tightlist
\item
  (Vidakis, 2024) Vidakis, N et al.~(2024). Reinforced HDPE with
  optimized biochar content for material extrusion additive
  manufacturing: morphological, rheological, electrical, and
  thermomechanical insights. Biochar, 6, 1-21.
  \url{https://doi.org/10.1007/s42773-024-00314-5}
\item
  (Bolanakis, 2024) Bolanakis, N et al.~(2024). Enhancing 3D Printing
  Materials with Biochar: A Literature Review. 2024 5th International
  Conference in Electronic Engineering, Information Technology \&
  Education (EEITE), 1-8.
  \url{https://doi.org/10.1109/EEITE61750.2024.10654408}
\item
  (Bolanakis, 2025) Bolanakis, N et al.~(2025). Valorization of Biochar
  as a Reinforcement Agent in Polyethylene Terephthalate Glycol for
  Additive Manufacturing: A Comprehensive Content Optimization Course.
  Journal of Manufacturing and Materials Processing.
  \url{https://doi.org/10.3390/jmmp9020068}
\item
  (Vidakis, 2023) Vidakis, N et al.~(2023). Biochar filler in MEX and
  VPP additive manufacturing: characterization and reinforcement effects
  in polylactic acid and standard grade resin matrices. Biochar, 5,
  1-21. \url{https://doi.org/10.1007/s42773-023-00238-6}
\item
  (Aboughaly, 2023) Aboughaly, M et al.~(2023). Enhancing the Potential
  of Polymer Composites Using Biochar as a Filler: A Review. Polymers,
  15. \url{https://doi.org/10.3390/polym15193981}
\item
  (George, 2023) George, J et al.~(2023). Improvement of Electrical and
  Mechanical Properties of PLA/PBAT Composites Using Coconut Shell
  Biochar for Antistatic Applications. Applied Sciences.
  \url{https://doi.org/10.3390/app13020902}
\item
  (Gkiliopoulos, 2025) Gkiliopoulos, D et al.~(2025). Valorization of
  Tomato Stem Waste: Biochar as a Filler in Three-Dimensional Printed
  PLA Composites. Polymers, 17.
  \url{https://doi.org/10.3390/polym17192565}
\item
  (Wijitkosum, 2023) Wijitkosum, S (2023). Influence of Pyrolysis
  Temperature and Time on Biochar Properties and Its Potential for
  Climate Change Mitigation. Journal of Human, Earth, and Future.
  \url{https://doi.org/10.28991/hef-2023-04-04-07}
\item
  (Suman, 2017) Suman, S, Gautam, S (2017). Effect of pyrolysis time and
  temperature on the characterization of biochars derived from biomass.
  Energy Sources, Part A: Recovery, Utilization, and Environmental
  Effects, 39, 933 - 940.
  \url{https://doi.org/10.1080/15567036.2016.1276650}
\item
  (Fernandes, 2020) Fernandes, B C C et al.~(2020). Impact of Pyrolysis
  Temperature on the Properties of Eucalyptus Wood-Derived Biochar.
  Materials, 13. \url{https://doi.org/10.3390/ma13245841}
\item
  (Zhang, 2024) Zhang, H et al.~(2024). Roles of biochars' properties in
  their water-holding capacity and bound water evaporation: quantitative
  importance and controlling mechanism. Biochar, 6, 1-14.
  \url{https://doi.org/10.1007/s42773-024-00317-2}
\item
  (Venkatesh, 2022) Venkatesh, G et al.~(2022). Characterization of
  Biochar Derived from Crop Residues for Soil Amendment, Carbon
  Sequestration and Energy Use. Sustainability.
  \url{https://doi.org/10.3390/su14042295}
\item
  (Hasan, 2019) Hasan, M M et al.~(2019). Effect of Pyrolysis
  Temperature and Time on Properties of Palm Kernel Shell-Based Biochar.
  IOP Conference Series: Materials Science and Engineering, 548.
  \url{https://doi.org/10.1088/1757-899X/548/1/012020}
\end{itemize}

\section{Chapter 2: Database
creation}\label{chapter-2-database-creation}

The data collected in this study include the pyrolysis conditions,
biochar properties and the properties of biochar composites produced
thorough additive manufacturing.

li \#\# 2.1 load packages

install.packages(c(``rmarkdown'',``knitr'',``tidyverse'',``DBI'',``RSQLite'',``dbplyr'',``kableExtra'',``ggpubr'',``here''))

library(DBI)

Lets establish a connection to the database:

\begin{Shaded}
\begin{Highlighting}[]
\NormalTok{biochar\_db }\OtherTok{\textless{}{-}} \FunctionTok{dbConnect}\NormalTok{(}\AttributeTok{drv =} \FunctionTok{SQLite}\NormalTok{(),}
                        \StringTok{"C:/Users/Owner/OneDrive {-} University of Idaho/Documents/Effect of pyrolysis condition on properties of biochar for additive manufacturing application/biochar.db"}\NormalTok{)}
\end{Highlighting}
\end{Shaded}

Load tables

\begin{Shaded}
\begin{Highlighting}[]
\FunctionTok{dbListTables}\NormalTok{(biochar\_db)}
\end{Highlighting}
\end{Shaded}

\begin{verbatim}
## [1] "BiocharPropertyValue" "BiocharSample"        "PropertyType"        
## [4] "PyrolysisCondition"   "Sheet1"
\end{verbatim}

\begin{Shaded}
\begin{Highlighting}[]
\NormalTok{biochar }\OtherTok{\textless{}{-}} \FunctionTok{dbGetQuery}\NormalTok{(biochar\_db, }\StringTok{"SELECT *FROM sheet1;"}\NormalTok{)}
\FunctionTok{head}\NormalTok{(biochar)}
\end{Highlighting}
\end{Shaded}

\begin{verbatim}
##   Sample_code pyrolysis_type temperature_(°C) carbon_content_(%)
## 1  BC_S450_r1           slow              450                 74
## 2  BC_S450_r2           slow              450                 74
## 3  BC_S450_r3           slow              450                 74
## 4  BC_S525_r1           slow              525                 78
## 5  BC_S525_r2           slow              525                 78
## 6  BC_S525_r3           slow              525                 78
##   nitrogen_content_(%) C/N_ratio density_g/cm3 raman_factor Flexural strength
## 1                 0.40  185.0000       0.52520     1.536032              2.83
## 2                 0.40  185.0000       0.52894     1.452542              2.65
## 3                 0.40  185.0000       0.53038     1.482350              2.99
## 4                 0.44  177.2727       0.55130     1.434615              2.56
## 5                 0.44  177.2727       0.52520     1.483202              2.75
## 6                 0.44  177.2727       0.52894     1.315257              2.58
##   Flexural modulus
## 1              926
## 2             1091
## 3             1036
## 4              468
## 5              510
## 6              470
\end{verbatim}

\begin{Shaded}
\begin{Highlighting}[]
\NormalTok{knitr}\SpecialCharTok{::}\FunctionTok{include\_graphics}\NormalTok{(}\StringTok{"C:/Users/Owner/OneDrive {-} University of Idaho/Pictures/Screenshots/Screenshot 2025{-}09{-}24 103730.png"}\NormalTok{)}
\end{Highlighting}
\end{Shaded}

\pandocbounded{\includegraphics[keepaspectratio]{../../Pictures/Screenshots/Screenshot 2025-09-24 103730.png}}
\#\# 2.1 Create table for pyrolysis condition

\begin{Shaded}
\begin{Highlighting}[]
\FunctionTok{dbExecute}\NormalTok{(biochar\_db, }\StringTok{"}
\StringTok{CREATE TABLE IF NOT EXISTS PyrolysisCondition (}
\StringTok{condition\_id INTEGER PRIMARY KEY,}
\StringTok{pyrolysis\_type TEXT NOT NULL,}
\StringTok{temperature\_c INTEGER NOT NULL,}
\StringTok{notes TEXT}
\StringTok{);}
\StringTok{"}\NormalTok{)}
\end{Highlighting}
\end{Shaded}

\begin{verbatim}
## [1] 0
\end{verbatim}

\subsection{2.2 Create Table Biochar
sample}\label{create-table-biochar-sample}

\begin{Shaded}
\begin{Highlighting}[]
\FunctionTok{dbExecute}\NormalTok{(biochar\_db, }\StringTok{"}
\StringTok{CREATE TABLE IF NOT EXISTS BiocharSample (}
\StringTok{sample\_id     INTEGER PRIMARY KEY,}
\StringTok{condition\_id  INTEGER NOT NULL,}
\StringTok{sample\_code   TEXT    NOT NULL UNIQUE,}
\StringTok{date\_produced TEXT,}
\StringTok{FOREIGN KEY (condition\_id) REFERENCES PyrolysisCondition(condition\_id) ON DELETE RESTRICT ON UPDATE CASCADE}
\StringTok{);}
\StringTok{"}\NormalTok{)}
\end{Highlighting}
\end{Shaded}

\begin{verbatim}
## [1] 0
\end{verbatim}

\subsection{2.3 CREATE table for Property
type}\label{create-table-for-property-type}

\begin{Shaded}
\begin{Highlighting}[]
\FunctionTok{dbExecute}\NormalTok{(biochar\_db, }\StringTok{"}
\StringTok{CREATE TABLE IF NOT EXISTS PropertyType (}
\StringTok{property\_id   INTEGER PRIMARY KEY,}
\StringTok{property\_name TEXT    NOT NULL UNIQUE,}
\StringTok{unit          TEXT,}
\StringTok{description   TEXT}
\StringTok{);}
\StringTok{"}\NormalTok{)}
\end{Highlighting}
\end{Shaded}

\begin{verbatim}
## [1] 0
\end{verbatim}

\subsection{2.4 CREATE TABLE FOR Biocharproperty
value}\label{create-table-for-biocharproperty-value}

\begin{Shaded}
\begin{Highlighting}[]
\FunctionTok{dbExecute}\NormalTok{(biochar\_db, }\StringTok{"}
\StringTok{CREATE TABLE IF NOT EXISTS BiocharPropertyValue (}
\StringTok{property\_value\_id INTEGER PRIMARY KEY,}
\StringTok{sample\_id         INTEGER NOT NULL,}
\StringTok{property\_id       INTEGER NOT NULL,}
\StringTok{value             TEXT    NOT NULL,}
\StringTok{method            TEXT,}
\StringTok{notes             TEXT,}
\StringTok{recorded\_at       TEXT DEFAULT (datetime(\textquotesingle{}now\textquotesingle{})),}
\StringTok{FOREIGN KEY (sample\_id)   REFERENCES BiocharSample(sample\_id) ON DELETE CASCADE ON UPDATE CASCADE,}
\StringTok{FOREIGN KEY (property\_id) REFERENCES PropertyType(property\_id) ON DELETE RESTRICT ON UPDATE CASCADE}
\StringTok{);}
\StringTok{"}\NormalTok{)}
\end{Highlighting}
\end{Shaded}

\begin{verbatim}
## [1] 0
\end{verbatim}

\section{Chapter 3: Data wrangling}\label{chapter-3-data-wrangling}

\begin{Shaded}
\begin{Highlighting}[]
\CommentTok{\# Install packages if needed}
\NormalTok{required\_packages }\OtherTok{\textless{}{-}} \FunctionTok{c}\NormalTok{(}\StringTok{"tidyverse"}\NormalTok{, }\StringTok{"DBI"}\NormalTok{, }\StringTok{"RSQLite"}\NormalTok{)}

\ControlFlowTok{for}\NormalTok{ (p }\ControlFlowTok{in}\NormalTok{ required\_packages) \{}
  \ControlFlowTok{if}\NormalTok{ (}\SpecialCharTok{!}\FunctionTok{require}\NormalTok{(p, }\AttributeTok{character.only =} \ConstantTok{TRUE}\NormalTok{)) \{}
    \FunctionTok{install.packages}\NormalTok{(p)}
    \FunctionTok{library}\NormalTok{(p, }\AttributeTok{character.only =} \ConstantTok{TRUE}\NormalTok{)}
\NormalTok{  \}}
\NormalTok{\}}

\CommentTok{\# Connect to the biochar database}
\NormalTok{biochar\_db }\OtherTok{\textless{}{-}} \FunctionTok{dbConnect}\NormalTok{(}
  \AttributeTok{drv =} \FunctionTok{SQLite}\NormalTok{(),}
  \AttributeTok{dbname =} \StringTok{"C:/Users/Owner/OneDrive {-} University of Idaho/Documents/Effect of pyrolysis condition on properties of biochar for additive manufacturing application/biochar.db"}
\NormalTok{)}

\CommentTok{\# Check available tables}
\FunctionTok{dbListTables}\NormalTok{(biochar\_db)}
\end{Highlighting}
\end{Shaded}

\begin{verbatim}
## [1] "BiocharPropertyValue" "BiocharSample"        "PropertyType"        
## [4] "PyrolysisCondition"   "Sheet1"
\end{verbatim}

\subsection{Load table from biochar}\label{load-table-from-biochar}

\begin{Shaded}
\begin{Highlighting}[]
\DocumentationTok{\#\# 🔄 **LOAD TABLES FROM \textasciigrave{}biochar.db\textasciigrave{}**  }

\CommentTok{\# Replace these with your actual table names}
\NormalTok{biochar\_raw }\OtherTok{\textless{}{-}} \FunctionTok{dbGetQuery}\NormalTok{(biochar\_db, }\StringTok{"SELECT * FROM Sheet1;"}\NormalTok{)}

\CommentTok{\# If you create additional tables later, load them here:}
\CommentTok{\# pyrolysis \textless{}{-} dbGetQuery(biochar\_db, "SELECT * FROM PyrolysisCondition;")}
\CommentTok{\# feedstock \textless{}{-} dbGetQuery(biochar\_db, "SELECT * FROM Feedstock;")}
\end{Highlighting}
\end{Shaded}

\subsection{Subsetting columns}\label{subsetting-columns}

\begin{Shaded}
\begin{Highlighting}[]
\DocumentationTok{\#\# 🧱 **SUBSETTING COLUMNS**}

\CommentTok{\# Example: select key variables (edit names to match your dataset)}
\NormalTok{biochar\_raw }\SpecialCharTok{\%\textgreater{}\%} 
  \FunctionTok{select}\NormalTok{(}\DecValTok{1}\SpecialCharTok{:}\DecValTok{5}\NormalTok{)   }\CommentTok{\# or select(colA, colB, colC)}
\end{Highlighting}
\end{Shaded}

\begin{verbatim}
##    Sample_code pyrolysis_type temperature_(°C) carbon_content_(%)
## 1   BC_S450_r1           slow              450                 74
## 2   BC_S450_r2           slow              450                 74
## 3   BC_S450_r3           slow              450                 74
## 4   BC_S525_r1           slow              525                 78
## 5   BC_S525_r2           slow              525                 78
## 6   BC_S525_r3           slow              525                 78
## 7   BC_S600_r1           slow              600                 84
## 8   BC_S600_r2           slow              600                 84
## 9   BC_S600_r3           slow              600                 84
## 10  BC_F450_r1           fast              450                 70
## 11  BC_F450_r2           fast              450                 70
## 12  BC_F450_r3           fast              450                 70
## 13  BC_F525_r1           fast              525                 77
## 14  BC_F525_r2           fast              525                 77
## 15  BC_F525_r3           fast              525                 77
## 16  BC_F600_r1           fast              600                 84
## 17  BC_F600_r2           fast              600                 84
## 18  BC_F600_r3           fast              600                 84
##    nitrogen_content_(%)
## 1                  0.40
## 2                  0.40
## 3                  0.40
## 4                  0.44
## 5                  0.44
## 6                  0.44
## 7                  0.51
## 8                  0.51
## 9                  0.51
## 10                 0.26
## 11                 0.26
## 12                 0.26
## 13                 0.41
## 14                 0.41
## 15                 0.41
## 16                 0.56
## 17                 0.56
## 18                 0.56
\end{verbatim}

\subsection{Subsetting rows}\label{subsetting-rows}

\begin{Shaded}
\begin{Highlighting}[]
\DocumentationTok{\#\# 🔍 **SUBSETTING ROWS**}
\CommentTok{\# Conditional filtering}
\NormalTok{biochar\_raw }\SpecialCharTok{\%\textgreater{}\%} \FunctionTok{filter}\NormalTok{(}\StringTok{\textasciigrave{}}\AttributeTok{temperature\_(°C)}\StringTok{\textasciigrave{}} \SpecialCharTok{\textgreater{}} \DecValTok{500}\NormalTok{)}
\end{Highlighting}
\end{Shaded}

\begin{verbatim}
##    Sample_code pyrolysis_type temperature_(°C) carbon_content_(%)
## 1   BC_S525_r1           slow              525                 78
## 2   BC_S525_r2           slow              525                 78
## 3   BC_S525_r3           slow              525                 78
## 4   BC_S600_r1           slow              600                 84
## 5   BC_S600_r2           slow              600                 84
## 6   BC_S600_r3           slow              600                 84
## 7   BC_F525_r1           fast              525                 77
## 8   BC_F525_r2           fast              525                 77
## 9   BC_F525_r3           fast              525                 77
## 10  BC_F600_r1           fast              600                 84
## 11  BC_F600_r2           fast              600                 84
## 12  BC_F600_r3           fast              600                 84
##    nitrogen_content_(%) C/N_ratio density_g/cm3 raman_factor Flexural strength
## 1                  0.44  177.2727       0.55130     1.434615             2.560
## 2                  0.44  177.2727       0.52520     1.483202             2.750
## 3                  0.44  177.2727       0.52894     1.315257             2.580
## 4                  0.51  164.7059       0.53038     1.380562             2.910
## 5                  0.51  164.7059       0.55150     1.395586             2.720
## 6                  0.51  164.7059       0.51140     1.444831             2.640
## 7                  0.41  187.8049       0.50850     1.474576             5.855
## 8                  0.41  187.8049       0.50872     1.477273             5.642
## 9                  0.41  187.8049       0.51512     1.474576             5.790
## 10                 0.56  150.0000       0.68180     1.292135             8.770
## 11                 0.56  150.0000       0.69200     1.316384             9.470
## 12                 0.56  150.0000       0.69140     1.188889             8.030
##    Flexural modulus
## 1               468
## 2               510
## 3               470
## 4               588
## 5               667
## 6               608
## 7              1780
## 8              1599
## 9              1575
## 10             3042
## 11             3087
## 12             3052
\end{verbatim}

\begin{Shaded}
\begin{Highlighting}[]
\CommentTok{\# Positional subsetting}
\NormalTok{biochar\_raw }\SpecialCharTok{\%\textgreater{}\%} \FunctionTok{slice}\NormalTok{(}\DecValTok{1}\SpecialCharTok{:}\DecValTok{10}\NormalTok{)}
\end{Highlighting}
\end{Shaded}

\begin{verbatim}
##    Sample_code pyrolysis_type temperature_(°C) carbon_content_(%)
## 1   BC_S450_r1           slow              450                 74
## 2   BC_S450_r2           slow              450                 74
## 3   BC_S450_r3           slow              450                 74
## 4   BC_S525_r1           slow              525                 78
## 5   BC_S525_r2           slow              525                 78
## 6   BC_S525_r3           slow              525                 78
## 7   BC_S600_r1           slow              600                 84
## 8   BC_S600_r2           slow              600                 84
## 9   BC_S600_r3           slow              600                 84
## 10  BC_F450_r1           fast              450                 70
##    nitrogen_content_(%) C/N_ratio density_g/cm3 raman_factor Flexural strength
## 1                  0.40  185.0000      0.525200     1.536032              2.83
## 2                  0.40  185.0000      0.528940     1.452542              2.65
## 3                  0.40  185.0000      0.530380     1.482350              2.99
## 4                  0.44  177.2727      0.551300     1.434615              2.56
## 5                  0.44  177.2727      0.525200     1.483202              2.75
## 6                  0.44  177.2727      0.528940     1.315257              2.58
## 7                  0.51  164.7059      0.530380     1.380562              2.91
## 8                  0.51  164.7059      0.551500     1.395586              2.72
## 9                  0.51  164.7059      0.511400     1.444831              2.64
## 10                 0.26  269.2308      0.381898     1.492308              2.04
##    Flexural modulus
## 1               926
## 2              1091
## 3              1036
## 4               468
## 5               510
## 6               470
## 7               588
## 8               667
## 9               608
## 10              611
\end{verbatim}

\subsection{Pipe operator}\label{pipe-operator}

\begin{Shaded}
\begin{Highlighting}[]
\DocumentationTok{\#\# 🧵 **PIPE OPERATOR EXAMPLE**}
\NormalTok{biochar\_filtered }\OtherTok{\textless{}{-}}\NormalTok{ biochar\_raw }\SpecialCharTok{\%\textgreater{}\%} 
  \FunctionTok{select}\NormalTok{(}\StringTok{\textasciigrave{}}\AttributeTok{temperature\_(°C)}\StringTok{\textasciigrave{}}\NormalTok{, }\StringTok{\textasciigrave{}}\AttributeTok{Flexural strength}\StringTok{\textasciigrave{}}\NormalTok{, }\StringTok{\textasciigrave{}}\AttributeTok{Flexural modulus}\StringTok{\textasciigrave{}}\NormalTok{) }\SpecialCharTok{\%\textgreater{}\%} 
  \FunctionTok{filter}\NormalTok{(}\StringTok{\textasciigrave{}}\AttributeTok{Flexural modulus}\StringTok{\textasciigrave{}} \SpecialCharTok{\textgreater{}} \DecValTok{10}\NormalTok{) }\SpecialCharTok{\%\textgreater{}\%} 
  \FunctionTok{slice}\NormalTok{(}\DecValTok{1}\NormalTok{)}
\NormalTok{biochar\_filtered}
\end{Highlighting}
\end{Shaded}

\begin{verbatim}
##   temperature_(°C) Flexural strength Flexural modulus
## 1              450              2.83              926
\end{verbatim}

\subsection{Converting temperature to
K}\label{converting-temperature-to-k}

\begin{Shaded}
\begin{Highlighting}[]
\DocumentationTok{\#\# ➕ **CREATING NEW COLUMNS**}
\NormalTok{biochar\_mutated }\OtherTok{\textless{}{-}}\NormalTok{ biochar\_raw }\SpecialCharTok{\%\textgreater{}\%} 
  \FunctionTok{mutate}\NormalTok{(}\AttributeTok{Temp\_K =} \StringTok{\textasciigrave{}}\AttributeTok{temperature\_(°C)}\StringTok{\textasciigrave{}} \SpecialCharTok{+} \FloatTok{273.15}\NormalTok{)}

\NormalTok{biochar\_mutated}
\end{Highlighting}
\end{Shaded}

\begin{verbatim}
##    Sample_code pyrolysis_type temperature_(°C) carbon_content_(%)
## 1   BC_S450_r1           slow              450                 74
## 2   BC_S450_r2           slow              450                 74
## 3   BC_S450_r3           slow              450                 74
## 4   BC_S525_r1           slow              525                 78
## 5   BC_S525_r2           slow              525                 78
## 6   BC_S525_r3           slow              525                 78
## 7   BC_S600_r1           slow              600                 84
## 8   BC_S600_r2           slow              600                 84
## 9   BC_S600_r3           slow              600                 84
## 10  BC_F450_r1           fast              450                 70
## 11  BC_F450_r2           fast              450                 70
## 12  BC_F450_r3           fast              450                 70
## 13  BC_F525_r1           fast              525                 77
## 14  BC_F525_r2           fast              525                 77
## 15  BC_F525_r3           fast              525                 77
## 16  BC_F600_r1           fast              600                 84
## 17  BC_F600_r2           fast              600                 84
## 18  BC_F600_r3           fast              600                 84
##    nitrogen_content_(%) C/N_ratio density_g/cm3 raman_factor Flexural strength
## 1                  0.40  185.0000      0.525200     1.536032             2.830
## 2                  0.40  185.0000      0.528940     1.452542             2.650
## 3                  0.40  185.0000      0.530380     1.482350             2.990
## 4                  0.44  177.2727      0.551300     1.434615             2.560
## 5                  0.44  177.2727      0.525200     1.483202             2.750
## 6                  0.44  177.2727      0.528940     1.315257             2.580
## 7                  0.51  164.7059      0.530380     1.380562             2.910
## 8                  0.51  164.7059      0.551500     1.395586             2.720
## 9                  0.51  164.7059      0.511400     1.444831             2.640
## 10                 0.26  269.2308      0.381898     1.492308             2.040
## 11                 0.26  269.2308      0.382670     1.492308             1.810
## 12                 0.26  269.2308      0.386570     1.492308             1.980
## 13                 0.41  187.8049      0.508500     1.474576             5.855
## 14                 0.41  187.8049      0.508720     1.477273             5.642
## 15                 0.41  187.8049      0.515120     1.474576             5.790
## 16                 0.56  150.0000      0.681800     1.292135             8.770
## 17                 0.56  150.0000      0.692000     1.316384             9.470
## 18                 0.56  150.0000      0.691400     1.188889             8.030
##    Flexural modulus Temp_K
## 1               926 723.15
## 2              1091 723.15
## 3              1036 723.15
## 4               468 798.15
## 5               510 798.15
## 6               470 798.15
## 7               588 873.15
## 8               667 873.15
## 9               608 873.15
## 10              611 723.15
## 11              533 723.15
## 12              580 723.15
## 13             1780 798.15
## 14             1599 798.15
## 15             1575 798.15
## 16             3042 873.15
## 17             3087 873.15
## 18             3052 873.15
\end{verbatim}

\subsection{Tibble}\label{tibble}

\begin{Shaded}
\begin{Highlighting}[]
\DocumentationTok{\#\# 📘 **CONVERT TO TIBBLE**}
\NormalTok{biochar\_tbl }\OtherTok{\textless{}{-}}\NormalTok{ biochar\_raw }\SpecialCharTok{\%\textgreater{}\%} \FunctionTok{as\_tibble}\NormalTok{()}
\NormalTok{biochar\_tbl}
\end{Highlighting}
\end{Shaded}

\begin{verbatim}
## # A tibble: 18 x 10
##    Sample_code pyrolysis_type `temperature_(°C)` `carbon_content_(%)`
##    <chr>       <chr>                       <int>                <int>
##  1 BC_S450_r1  slow                          450                   74
##  2 BC_S450_r2  slow                          450                   74
##  3 BC_S450_r3  slow                          450                   74
##  4 BC_S525_r1  slow                          525                   78
##  5 BC_S525_r2  slow                          525                   78
##  6 BC_S525_r3  slow                          525                   78
##  7 BC_S600_r1  slow                          600                   84
##  8 BC_S600_r2  slow                          600                   84
##  9 BC_S600_r3  slow                          600                   84
## 10 BC_F450_r1  fast                          450                   70
## 11 BC_F450_r2  fast                          450                   70
## 12 BC_F450_r3  fast                          450                   70
## 13 BC_F525_r1  fast                          525                   77
## 14 BC_F525_r2  fast                          525                   77
## 15 BC_F525_r3  fast                          525                   77
## 16 BC_F600_r1  fast                          600                   84
## 17 BC_F600_r2  fast                          600                   84
## 18 BC_F600_r3  fast                          600                   84
## # i 6 more variables: `nitrogen_content_(%)` <dbl>, `C/N_ratio` <dbl>,
## #   `density_g/cm3` <dbl>, raman_factor <dbl>, `Flexural strength` <dbl>,
## #   `Flexural modulus` <int>
\end{verbatim}

\subsection{Joining tables}\label{joining-tables}

\begin{Shaded}
\begin{Highlighting}[]
\DocumentationTok{\#\# 🔗 **JOINING TABLES**  }
\CommentTok{\# Example join (uncomment when you have second table)}
\CommentTok{\# biochar\_joined \textless{}{-} biochar\_raw \%\textgreater{}\% }
\CommentTok{\#   left\_join(pyrolysis, by = "condition\_id")}
\CommentTok{\# biochar\_joined}
\end{Highlighting}
\end{Shaded}

\subsection{Groupwise summary}\label{groupwise-summary}

\begin{Shaded}
\begin{Highlighting}[]
\DocumentationTok{\#\# 📊 **GROUPWISE SUMMARIES**}
\CommentTok{\# Min, mean, max Flexural strength}
\NormalTok{biochar\_raw }\SpecialCharTok{\%\textgreater{}\%} 
  \FunctionTok{group\_by}\NormalTok{(pyrolysis\_type) }\SpecialCharTok{\%\textgreater{}\%} 
  \FunctionTok{summarize}\NormalTok{(}
    \AttributeTok{min\_fs =} \FunctionTok{min}\NormalTok{(}\StringTok{\textasciigrave{}}\AttributeTok{Flexural strength}\StringTok{\textasciigrave{}}\NormalTok{, }\AttributeTok{na.rm =} \ConstantTok{TRUE}\NormalTok{),}
    \AttributeTok{mean\_fs =} \FunctionTok{mean}\NormalTok{(}\StringTok{\textasciigrave{}}\AttributeTok{Flexural strength}\StringTok{\textasciigrave{}}\NormalTok{, }\AttributeTok{na.rm =} \ConstantTok{TRUE}\NormalTok{),}
    \AttributeTok{max\_fs =} \FunctionTok{max}\NormalTok{(}\StringTok{\textasciigrave{}}\AttributeTok{Flexural strength}\StringTok{\textasciigrave{}}\NormalTok{, }\AttributeTok{na.rm =} \ConstantTok{TRUE}\NormalTok{)}
\NormalTok{  ) }\SpecialCharTok{\%\textgreater{}\%} 
  \FunctionTok{arrange}\NormalTok{(}\FunctionTok{desc}\NormalTok{(mean\_fs))}
\end{Highlighting}
\end{Shaded}

\begin{verbatim}
## # A tibble: 2 x 4
##   pyrolysis_type min_fs mean_fs max_fs
##   <chr>           <dbl>   <dbl>  <dbl>
## 1 fast             1.81    5.49   9.47
## 2 slow             2.56    2.74   2.99
\end{verbatim}

\begin{Shaded}
\begin{Highlighting}[]
\CommentTok{\# Min, mean, max Flexural modulus}
\NormalTok{biochar\_raw }\SpecialCharTok{\%\textgreater{}\%} 
  \FunctionTok{group\_by}\NormalTok{(pyrolysis\_type) }\SpecialCharTok{\%\textgreater{}\%} 
  \FunctionTok{summarize}\NormalTok{(}
    \AttributeTok{min\_fm =} \FunctionTok{min}\NormalTok{(}\StringTok{\textasciigrave{}}\AttributeTok{Flexural modulus}\StringTok{\textasciigrave{}}\NormalTok{, }\AttributeTok{na.rm =} \ConstantTok{TRUE}\NormalTok{),}
    \AttributeTok{mean\_fm =} \FunctionTok{mean}\NormalTok{(}\StringTok{\textasciigrave{}}\AttributeTok{Flexural modulus}\StringTok{\textasciigrave{}}\NormalTok{, }\AttributeTok{na.rm =} \ConstantTok{TRUE}\NormalTok{),}
    \AttributeTok{max\_fm =} \FunctionTok{max}\NormalTok{(}\StringTok{\textasciigrave{}}\AttributeTok{Flexural modulus}\StringTok{\textasciigrave{}}\NormalTok{, }\AttributeTok{na.rm =} \ConstantTok{TRUE}\NormalTok{)}
\NormalTok{  ) }\SpecialCharTok{\%\textgreater{}\%} 
  \FunctionTok{arrange}\NormalTok{(}\FunctionTok{desc}\NormalTok{(mean\_fm))}
\end{Highlighting}
\end{Shaded}

\begin{verbatim}
## # A tibble: 2 x 4
##   pyrolysis_type min_fm mean_fm max_fm
##   <chr>           <int>   <dbl>  <int>
## 1 fast              533   1762.   3087
## 2 slow              468    707.   1091
\end{verbatim}

\subsection{Extracting a single
column}\label{extracting-a-single-column}

\begin{Shaded}
\begin{Highlighting}[]
\NormalTok{temps }\OtherTok{\textless{}{-}}\NormalTok{ biochar\_raw }\SpecialCharTok{\%\textgreater{}\%} \FunctionTok{pull}\NormalTok{(}\StringTok{\textasciigrave{}}\AttributeTok{temperature\_(°C)}\StringTok{\textasciigrave{}}\NormalTok{)}
\NormalTok{temps}
\end{Highlighting}
\end{Shaded}

\begin{verbatim}
##  [1] 450 450 450 525 525 525 600 600 600 450 450 450 525 525 525 600 600 600
\end{verbatim}

\subsection{Conditional value
assignment}\label{conditional-value-assignment}

\begin{verbatim}
##   carbon_content_(%) n
## 1               High 6
## 2                Low 6
## 3             Medium 6
\end{verbatim}

\section{Chapter 4: Data
visualization}\label{chapter-4-data-visualization}

\subsection{4.1 Temperature vs carbon
content}\label{temperature-vs-carbon-content}

\begin{Shaded}
\begin{Highlighting}[]
\CommentTok{\# Scatterplot: Temperature vs carbon\_content**}
\FunctionTok{ggplot}\NormalTok{(biochar, }\FunctionTok{aes}\NormalTok{(}\AttributeTok{x =} \StringTok{\textasciigrave{}}\AttributeTok{temperature\_(°C)}\StringTok{\textasciigrave{}}\NormalTok{, }\AttributeTok{y =} \StringTok{\textasciigrave{}}\AttributeTok{carbon\_content\_(\%)}\StringTok{\textasciigrave{}}\NormalTok{)) }\SpecialCharTok{+}
  \FunctionTok{geom\_point}\NormalTok{(}\AttributeTok{alpha =} \FloatTok{0.7}\NormalTok{) }\SpecialCharTok{+}
  \FunctionTok{geom\_smooth}\NormalTok{(}\AttributeTok{method =} \StringTok{"lm"}\NormalTok{, }\AttributeTok{se =} \ConstantTok{FALSE}\NormalTok{, }\AttributeTok{linewidth =} \DecValTok{1}\NormalTok{) }\SpecialCharTok{+}
  \FunctionTok{theme\_minimal}\NormalTok{() }\SpecialCharTok{+}
  \FunctionTok{labs}\NormalTok{(}
    \AttributeTok{title =} \StringTok{"Relationship Between Temperature and carbon content"}\NormalTok{,}
    \AttributeTok{x =} \StringTok{"Temperature (°C)"}\NormalTok{,}
    \AttributeTok{y =} \StringTok{"Carbon content"}
\NormalTok{  )}
\end{Highlighting}
\end{Shaded}

\begin{verbatim}
## `geom_smooth()` using formula = 'y ~ x'
\end{verbatim}

\pandocbounded{\includegraphics[keepaspectratio]{biochar_files/figure-latex/scatterplot of temp against carbon content-1.pdf}}
\#\# 4.2 Box plot

\begin{Shaded}
\begin{Highlighting}[]
\FunctionTok{ggplot}\NormalTok{(biochar, }\FunctionTok{aes}\NormalTok{(}\AttributeTok{x =}\NormalTok{ pyrolysis\_type, }\AttributeTok{y =} \StringTok{\textasciigrave{}}\AttributeTok{Flexural modulus}\StringTok{\textasciigrave{}}\NormalTok{)) }\SpecialCharTok{+}
  \FunctionTok{geom\_boxplot}\NormalTok{(}\AttributeTok{fill =} \StringTok{"skyblue"}\NormalTok{) }\SpecialCharTok{+}
  \FunctionTok{theme\_minimal}\NormalTok{() }\SpecialCharTok{+}
  \FunctionTok{labs}\NormalTok{(}
    \AttributeTok{title =} \StringTok{"Effect of pyrolysis type vs flexural strength"}\NormalTok{,}
    \AttributeTok{x =} \StringTok{"pyrolysis type"}\NormalTok{,}
    \AttributeTok{y =} \StringTok{"Flexural modulus"}
\NormalTok{  ) }\SpecialCharTok{+}
  \FunctionTok{theme}\NormalTok{(}\AttributeTok{axis.text.x =} \FunctionTok{element\_text}\NormalTok{(}\AttributeTok{angle =} \DecValTok{45}\NormalTok{, }\AttributeTok{hjust =} \DecValTok{1}\NormalTok{))}
\end{Highlighting}
\end{Shaded}

\pandocbounded{\includegraphics[keepaspectratio]{biochar_files/figure-latex/box plot-1.pdf}}

\end{document}
